\section{The cotangent envelope: compactness and frames}
\label{sec:discs}

We study the flow through boundary averages of its dual squared norm on
holomorphic cotangent discs.  Let $\Delta=\{z\in\C:|z|<1\}$, let $dA$
denote Euclidean area on $\Delta$, and write
$\pi:Y=T^{*(1,0)}M\to M$.  Put
\[
 N(x,\xi,t)=|\xi|^2_{h_t^{-1}}.
\]
Thus $N(\cdot,t)\ge N(\cdot,0)$ by
Theorem~\ref{thm:localized-flow}.  The function $N(\cdot,0)$ is
plurisubharmonic: the squared norm on the total space of a Griffiths
nonpositive bundle is plurisubharmonic, and the dual of $T^{1,0}M$ has
this curvature sign.

Fix $T<\infty$ and denote physical time at a disc center by $s$.
If $y\in Y$ and $0<s<T$, let $v^{\rm raw}(y,s)$ be the infimum of
\begin{equation}
 \left\langle N\bigl(f(\zeta),s e^{w(\zeta)+\overline{w(\zeta)}}\bigr)
 \right\rangle,
 \qquad (f(0),w(0))=(y,0),
 \label{var:envelope}
\end{equation}
over holomorphic discs $f:\overline\Delta\to Y$ and
$w:\overline\Delta\to\C$, smooth on the closed unit disc, with
$s e^{w+\bar w}<T$ there.  Brackets denote normalized circle average.
Let $v$ be the lower semicontinuous regularization in $(y,s)$.
Constant discs, the plurisubharmonicity of $N(\cdot,0)$, and
$N(\cdot,t)\ge N(\cdot,0)$ give
\begin{equation}
 N(y,0)\le v(y,s)\le v^{\rm raw}(y,s)\le N(y,s).
 \label{var:sandwich}
\end{equation}
Multiplying the fiber component of every competing disc by a nonzero
complex number shows that $v$ is homogeneous of degree two in the fiber.
The assertion at the zero section follows from \eqref{var:sandwich}.

Minimizing boundary averages over discs with a prescribed center is the
Poisson-envelope construction in the theory of Poletsky and Rosay
\cite{Poletsky1991,Rosay2003}; see also
\cite[Sections~1--2]{LarussonSigurdsson} for the manifold framework.
Here we prove the differential inequality needed for the evolving norm
directly, with the compactness and quantitative deformation estimates
specified below.

Our goal is to prove $v=N$.  This equality will give the submean
inequality for $N(x,\xi,e^{w+\bar w})$, hence its joint
plurisubharmonicity.  The competitors in \eqref{var:envelope} have no
common boundary compact or area bound.  In Section~\ref{sec:variation}
we will localize almost minimizing discs by boundary walls and a positive
area penalty.  The compact containing their full images may then depend
on the penalty.  With centers in a fixed compact coordinate region, the
frame estimates used in the differential inequality must depend on
boundary and center data, independently of the area bound and the
compact containing the interiors.  This Section establishes these two
kinds of control, together with the cotangent evolution identity used
in that estimate.  All compactness arguments use the complete initial
metric $g$.

\subsection{The symplectic Hessian operator}

The canonical holomorphic bivector on $Y$ is
\[
 \mathcal J=\sum_{i=1}^n
       \frac{\partial}{\partial x_i}\wedge
       \frac{\partial}{\partial\xi_i}.
\]
We use column matrices for Hermitian forms, so a form with matrix $W$
has quadratic value $v^*Wv$.  In a frame $e_1,\ldots,e_{2n}$, write
$\mathcal J=\sum_{A<B}J^{AB}e_A\wedge e_B$ and define
\begin{equation}\label{disc:p-definition}
 p(W,\mathcal J)=
 \sum_{A<B,\,C<D}\overline{J^{AB}}J^{CD}
       (W_{AC}W_{BD}-W_{AD}W_{BC}).
\end{equation}
This is the induced exterior-square squared norm when $W$ is positive,
and is a real polynomial for every Hermitian $W$.  The definition is
independent of frame.  In canonical cotangent coordinates we abbreviate
it to $p(W)$; if
$W=\left(\begin{smallmatrix}A&B\\B^*&C\end{smallmatrix}\right)$, then
\begin{equation}\label{disc:p-block}
 p(W)=\tr(AC^{\mathsf T})-
                 \sum_{i,j}B_{ij}\overline{B_{ji}}.
\end{equation}
Here and below the transpose records the tangent--cotangent pairing.
In particular, if tangent coefficients have metric matrix $h$, the
column matrix of the dual norm is $b=h^{-\mathsf T}$.

\begin{lemma}[Symplectic evolution]\label{lem:symplectic-evolution}
For the flow of Theorem~\ref{thm:localized-flow},
\begin{equation}\label{disc:symplectic-evolution}
 \partial_tN=p(\dd_YN).
\end{equation}
For every real smooth base function $d=d(x,t)$, with $t$ held fixed in
the spatial Hessian,
\begin{equation}\label{disc:base-addition}
 p(\dd_Y(N+d))=\partial_tN+\Delta_{h_t}d.
\end{equation}
These identities do not assume that $h_t$ is complete or has
nonnegative curvature at positive time.
\end{lemma}

\begin{proof}
Fix $(x,t)$ and use K\"ahler-normal coordinates for $h_t$ there.
At that point $h_t=b_t=I$, all first spatial derivatives of $b_t$
vanish, and the Hessian of $N$ has blocks
\[
 \dd_YN=\begin{pmatrix}\mathcal R_\xi&0\\0&I\end{pmatrix}.
\]
The horizontal form $\mathcal R_\xi$ is the bisectional curvature form
of $h_t$ with one pair of indices contracted against the metric dual
of $\xi$.  Indeed differentiating $b=h^{-\mathsf T}$ twice at a normal
point gives the negative second derivatives of $h$, transposed, which
are precisely those curvature components.  The K\"ahler symmetries give
\[
 \tr\mathcal R_\xi=\xi^*\Ric(h_t)^{\mathsf T}\xi.
\]
On the other hand $\partial_th=-\Ric(h)$ implies
$\partial_tb=h^{-\mathsf T}\Ric(h)^{\mathsf T}h^{-\mathsf T}$.
Equation~\eqref{disc:p-block} now proves
\eqref{disc:symplectic-evolution}.  Adding a base Hessian changes only
the horizontal block.  Its contribution to $p$ is
$\tr((\dd_xd)b^{\mathsf T})=\tr(h_t^{-1}\dd_xd)$, proving
\eqref{disc:base-addition}.  The calculation is tensorial, so the
normal-coordinate verification proves both identities everywhere.
\end{proof}

\subsection{Compact boundary and bounded area}

\begin{proposition}[Compactness of discs]\label{prop:disc-compactness}
For each compact $K\subset M$ and each $A_0<\infty$, there is a
compact $K'\subset M$ such that every holomorphic disc
$f:\overline\Delta\to M$, smooth up to the circle, satisfying
\[
 f(\partial\Delta)\subset K,\qquad
 \int_\Delta f^*\omega_g\le A_0
\]
has $f(\overline\Delta)\subset K'$.  Every sequence of such discs
has a subsequence converging holomorphically on compact subsets of
$\Delta$.

The same statements hold for discs in $Y$ with compact boundary in
$Y$ and with $\int_\Delta(\pi\circ f)^*\omega_g\le A_0$.
The resulting compact may depend on $K$ and $A_0$.
\end{proposition}

The area bound first places part of each disc in a fixed compact.
An inverse-Jacobian estimate for the global injective immersion then
propagates this control away from isolated possible singularities of a
limit; properness of the immersion is not assumed.  The following
extra-integrability lemma removes those punctures.  A second lemma
records the subharmonic compactness needed to propagate the initial
control.

The puncture estimate assumes no extension of the bundle or its metric
across the puncture.

\begin{lemma}[Extra integrability]\label{lem:extra-integrability}
Let $(E,k)$ be a smooth Hermitian holomorphic vector bundle on
$\Delta^*=\{0<|z|<1\}$, with Griffiths nonnegative curvature.
Suppose that $A:E\to\Delta^*\times\C^q$ is holomorphic, pointwise
injective, and has bounded operator norm.  If $a$ is a holomorphic
section with $\int_{\Delta^*}|a|_k^2\,dA<\infty$, then there are
$r_0>0$ and $\sigma>0$ such that
\begin{equation}\label{disc:extra-integrability}
 \int_{0<|z|<r_0}|a|_k^2|z|^{-\sigma}\,dA<\infty.
\end{equation}
The exponent may depend on the section and the bundle data.
\end{lemma}

\begin{proof}
Choose a fixed smooth function $\chi:\R\to[0,1]$ which is zero on
$(-\infty,-2]$ and one on $[-1,\infty)$, and set
\[
 \chi_m(z)=\chi(m^{-1}\log|z|),\qquad
 \phi_m(z)=\frac2m\log|z|+|z|^{1/m}.
\]
Writing $r=|z|$, one has
\[
 |\dbar\chi_m|\le\frac{C}{mr},\qquad
 (\phi_m)_{z\bar z}=\frac1{4m^2}r^{1/m-2}.
\]
On a complex curve Griffiths and Nakano nonnegativity coincide.  Apply
the bundle-valued H\"ormander estimate to $E$-valued $(1,0)$-forms,
with the metric $ke^{-\phi_m}$ and source
$a\,\dbar\chi_m\wedge dz$; see \cite{Hormander} and
\cite[Theorem 5.1 and Remark 4.5]{Demailly}.
One may use any complete auxiliary K\"ahler metric on $\Delta^*$.
In dimension one its conformal factor cancels both from the norm of
the unknown top form times volume and from the curvature-inverse
source pairing times volume.  Dropping the nonnegative curvature of
$k$ therefore gives a solution $v_m$ of
$\dbar v_m=a\,\dbar\chi_m$ satisfying
\begin{align}
 \int_{\Delta^*}|v_m|_k^2e^{-\phi_m}\,dA
 &\le C\int_{e^{-2m}<r<e^{-m}}
       |a|_k^2r^{-3/m}e^{-r^{1/m}}\,dA,\label{disc:weighted-solve}\\
 \int_{\Delta^*}|v_m|_k^2r^{-2/m}\,dA
 &\le C\int_{e^{-2m}<r<e^{-m}}|a|_k^2\,dA\longrightarrow0.
 \label{disc:correction-small}
\end{align}
For the second estimate we used $r^{-3/m}\le e^6$ on the source
annulus and $e^{-1}\le e^{-r^{1/m}}\le1$.  The constants are
independent of $m$.  The estimate concerns smooth metrics on the
punctured domain; it does not assert a curvature extension at zero.

The sections $a_m=\chi_ma-v_m$ are holomorphic and converge to $a$
in unweighted $L^2$.  Near zero they equal $-v_m$, so each has a
positive extra integrability exponent, namely $2/m$.  Since $A$ is
bounded, $Aa_m$ and $Aa$ are ordinary $L^2$ holomorphic vector
functions on the punctured disc.  Their negative Laurent coefficients
vanish by $L^2$ integrability; hence they extend across zero.
Their $L^2$ convergence and the interior Cauchy estimates give
convergence of every finite jet at zero.

Here is the closure argument that recovers a positive exponent for
$a$ itself.  Let $R=\mathcal O_{\C,0}$, with maximal ideal
$\mathfrak m=(z)$, and let $S\subset R^q$ consist of all germs $Ab$
where $b$ is holomorphic on some punctured neighborhood and
\[
 \int |b|_k^2r^{-\varepsilon}\,dA<\infty
 \quad\hbox{for some }\varepsilon>0.
\]
These images extend across zero by the same $L^2$ argument.  The set
$S$ is an $R$-submodule: for a sum take the smaller of the two positive
exponents, and multiplication by a holomorphic germ is bounded after
shrinking the neighborhood.  The ring $R$ is Noetherian, so $S$ is a
submodule of a finite free module and $Q=R^q/S$ is a finite module.
For each $k$, the image of $S$ in $R^q/\mathfrak m^kR^q$ is a
finite-dimensional vector subspace and is closed.  Jet convergence of
$Aa_m\in S$ consequently gives
\[
 Aa\in S+\mathfrak m^kR^q\qquad\hbox{for every }k.
\]
Krull's intersection Theorem for the finite local module $Q$ gives
$\bigcap_k\mathfrak m^kQ=0$ \cite[Lemma 10.51.4]{StacksKrull}.
Thus $Aa\in S$.  By definition it is the image of one section $b$
with some positive extra exponent.  Pointwise injectivity of $A$
implies $a=b$ on a sufficiently small punctured neighborhood, proving
\eqref{disc:extra-integrability}.  No common exponent for $a_m$ was
used.
\end{proof}

Let $F=(F_1,\ldots,F_N):M\to\C^N$ be the global injective
holomorphic immersion of Proposition~\ref{prop:initial-data}.  Set
\[
 \psi=\log\sum_{|I|=n}|dF_I|_g^2,
 \qquad dF_I=dF_{i_1}\wedge\cdots\wedge dF_{i_n}.
\]
Immersivity makes $\psi$ smooth.  Along any holomorphic disc,
$\psi$ is subharmonic: each $dF_I$ is a holomorphic canonical form,
whose logarithmic squared norm is plurisubharmonic under the curvature
hypothesis, and the log-sum has the same property.  Similarly
$|dF|_g^2=\sum_i|dF_i|_g^2$ is plurisubharmonic.

\begin{lemma}[Compactness with a lower anchor]\label{lem:psh-compactness-bridge}
Let $\Omega_j$ be increasing open subsets exhausting a connected
complex manifold $X$, and let $u_j$ be plurisubharmonic on $\Omega_j$.
Assume that for each compact $K\subset X$, the functions $u_j$ are
uniformly bounded above on $K$ once $K\subset\Omega_j$.
Every sequence has a subsequence for which either $u_j$ tends to
$-\infty$ locally uniformly from above, or $u_j$ converges in
$L^1_{\mathrm{loc}}$ to a plurisubharmonic function.
A lower bound $u_j(z_j)\ge-A$ at points $z_j$ in a fixed compact set
excludes the first alternative along that sequence.  In particular,
lower bounds on moving subsets of a fixed compact set of uniformly
positive measure suffice.  If the $L^1_{\mathrm{loc}}$ limit is the
constant $c$, then
\[
 \limsup_j\sup_K u_j\le c
\]
for every compact $K\subset X$.
\end{lemma}

\begin{proof}
We explain the lower-anchor step rather than building it into the
compactness theorem.  Pass to a subsequence with $z_j\to z_*$ and
work in a coordinate ball about $z_*$.  On a slightly larger ball
choose a common upper bound $C$ and set $w_j=C-u_j$.
These functions are nonnegative and superharmonic.  The supermean
inequality at $z_j$ bounds their integrals on a ball of fixed radius
about $z_j$ by its volume times $C+A$.  A fixed smaller ball about
$z_*$ is eventually contained in those balls, so the $L^1$ norms
there are bounded.

Such a bound propagates along finite chains of overlapping coordinate
balls.  In a fixed overlap of positive measure, the integral bound
supplies a point at which $w_j$ is bounded by its average over the
overlap.  Applying the supermean inequality on a larger ball centered
at that point, still contained in the coordinate neighborhood, bounds
the integral on the next smaller ball.  On each fixed compact overlap, the Euclidean volume measures in
the two coordinate charts are uniformly comparable.  Increasing the
local upper bound $C$ when moving between charts only adds a fixed constant.
Connectedness and a finite covering therefore give $L^1$ bounds on
every fixed compact subset of $X$.

Families of subharmonic functions bounded in $L^1_{\mathrm{loc}}$
are relatively compact in that topology \cite[Chapter~I, Proposition~4.21]{DemaillyBook}.
Apply this result in coordinate balls and take a diagonal subsequence.
Positivity of the complex Hessians passes to the distributional
limit, so its upper semicontinuous representative is
plurisubharmonic.  If no subsequence has a lower anchor in any fixed
compact set, then the suprema on each member of a compact exhaustion
tend to $-\infty$; this is the other alternative.  This also proves
the assertion for exhausting domains, since each finite chain of
balls is contained in every sufficiently late $\Omega_j$.

Finally suppose the limit is the constant $c$.  In a coordinate ball
whose concentric enlargement is relatively compact, apply the
submean inequality on a fixed-radius ball centered at each point of
the smaller ball.  It bounds $u_j-c$ from above by a fixed constant
times the $L^1$ norm of $u_j-c$ on the enlargement.  This norm tends
to zero.  A finite covering proves the asserted upper bound on $K$.
\end{proof}

\begin{proof}[Proof of Proposition~\ref{prop:disc-compactness}]
We first control centers of discs $f_l$ in $M$.  The maximum principle
gives uniform bounds
\begin{equation}\label{disc:boundary-upper}
 |F\circ f_l|\le C_K,\qquad
 |dF|_g\circ f_l\le C_K,\qquad
 \psi_l:=\psi\circ f_l\le C_K.
\end{equation}
The area bound controls the energy, with a fixed normalization
constant.  In particular, the functions
\[
 E_l(\theta)=\int_{1/2}^1
          |\partial_rf_l(re^{i\theta})|_g^2r\,dr
\]
have uniformly bounded angular integrals.  On a set of angles of
measure bounded below, $E_l(\theta)\le C(A_0+1)$.  For each such
angle and every $r\in[1/2,3/4]$, Cauchy--Schwarz gives
\[
 \dist_g(f_l(re^{i\theta}),K)
 \le E_l(\theta)^{1/2}
       \left(\int_{1/2}^1\frac{dr}{r}\right)^{1/2}.
\]
Completeness and Hopf--Rinow put these images in a fixed compact.
Thus on a subset of the fixed annulus of area bounded below,
$\psi_l$ has a fixed lower bound.  Lemma~\ref{lem:psh-compactness-bridge}, together with
\eqref{disc:boundary-upper}, now gives a subsequence converging in
$L^1_{\mathrm{loc}}(\Delta)$ to a subharmonic function
$\psi_\infty$; the lower bounds on the anchor subsets exclude
collapse to $-\infty$.  We retain their positive-area property for
the later step that avoids neighborhoods of the possible atoms.  Let
$\mu_l=(2\pi)^{-1}\Delta_{\mathrm{eucl}}\psi_l$ and
$\mu=(2\pi)^{-1}\Delta_{\mathrm{eucl}}\psi_\infty$.
Then $\mu_l\to\mu$ weakly on compact subsets.

We spell out the required inverse-Jacobian estimate.  Let
$Z=\{z:\mu(\{z\})\ge1\}$, a locally finite subset of $\Delta$.
Near any point outside $Z$, choose a small disc $D$ whose boundary
has zero $\mu$-measure and for which $\mu(D)<b<4/3$.
For large $l$, $\mu_l(D)<b$.  The Riesz decomposition on $D$
\cite[Chapter~I, Proposition~4.22]{DemaillyBook} is
\[
 \psi_l=h_l+\int_D\log|z-\zeta|\,d\mu_l(\zeta).
\]
On each smaller concentric disc the harmonic terms $h_l$ are bounded
in absolute value: their $L^1$ norms are bounded by $L^1$ convergence
of $\psi_l$, the bounded masses, and local integrability of the
logarithmic kernel; harmonic interior estimates then apply.
For a positive measure $\nu$ of mass $b_l\le b$, Jensen's inequality
gives, with the zero-mass case interpreted separately,
\[
 \exp\left(-p\int_D\log|z-\zeta|\,d\nu(\zeta)\right)
 \le\frac1{b_l}\int_D|z-\zeta|^{-pb_l}\,d\nu(\zeta),
 \qquad p=\tfrac32.
\]
Integrating in $z$ gives a uniform bound since $pb<2$.  Therefore
$e^{-\psi_l}$ is locally bounded in $L^{3/2}(\Delta\setminus Z)$.

Let $s_1\le\cdots\le s_n$ be the singular values of $dF$ measured
from $g$ to the Euclidean metric.  Cauchy--Binet and
\eqref{disc:boundary-upper} imply
\[
 e^{\psi_l}=\prod_i s_i(f_l)^2,\qquad
 s_1(f_l)^{-2}\le C_K e^{-\psi_l}.
\]
Consequently
\begin{equation}\label{disc:inverse-jacobian}
 |f_l'|_g^2\le C_K e^{-\psi_l}|(F\circ f_l)'|^2.
\end{equation}
The bounded holomorphic maps $F\circ f_l$ have uniformly bounded
derivatives on smaller discs.  Hence $df_l$ is locally bounded in
$L^3$ off $Z$.  For $p\in M$ the function
$v_{l,p}(z)=\dist_g(f_l(z),p)$ has weak gradient bounded by
$|df_l|_g$.  Morrey's local gradient-only estimate in real dimension
two, with exponent $1-2/3=1/3$, therefore applies
\cite[Section~5.6.2, Theorem~4 and the Remark on p.~283]{EvansPDE}.
Its constant depends on the nested domain discs, not on $p$;
no bound on the values of $v_{l,p}$ is needed.
Taking $p=f_l(w)$ gives the uniform local metric H\"older estimate.

These local estimates propagate the compact anchors.  Indeed, remove
small neighborhoods of the finitely many points of $Z$ in the anchor
annulus, with total area smaller than half the anchor lower bound.
Choose anchor points in the remaining compact set and pass to a
convergent subsequence of their domain positions.  The local
H\"older estimate puts their neighboring images in a fixed bounded
$g$-neighborhood of the anchor compact.  Any compact subset of the
connected set $\Delta\setminus Z$ can be reached by finitely many
overlapping small discs of this kind.  Completeness makes each
successive bounded neighborhood compact.  Arzel\`a--Ascoli, local
holomorphic coordinates, and diagonal extraction give a holomorphic
limit $f:\Delta\setminus Z\to M$.  Its area is finite by lower
semicontinuity, and $dF$ remains uniformly bounded along it.

Fix a puncture of this limit and move its domain coordinate to zero.
The bundle $E=f^*T^{1,0}M$ has Griffiths nonnegative curvature,
$A=dF:E\to\C^N$ is bounded and injective, and the holomorphic
section $a=f'$ has finite $L^2$ norm.  Lemma~\ref{lem:extra-integrability}
applies.  Fubini then gives an angle $\theta$ with
\[
 \int_0^{r_0}|a(re^{i\theta})|_g^2r^{1-\sigma}\,dr<\infty.
\]
Since $\int_0^{r_0}r^{\sigma-1}\,dr<\infty$, the corresponding
radial path has finite length and converges to a point $p\in M$.
The bounded map $F\circ f$ extends holomorphically to zero and its
value there is $F(p)$.  Choose a relatively compact coordinate
neighborhood $U$ of $p$ with compact boundary.  Injectivity of $F$
gives
\[
 \dist_{\mathrm{eucl}}(F(p),F(\partial U))>0.
\]
For small enough $r$, $F\circ f$ on $0<|z|<r$ avoids
$F(\partial U)$, and the radial path supplies a point with image in
$U$.  Connectedness traps the entire punctured image in $U$.
Bounded coordinate functions extend, so $f$ extends holomorphically.
This argument uses compactness of a local chart boundary and global
injectivity of $F$; it does not require $F$ to be proper.

Off the puncture, $\psi_\infty=\psi\circ f$ almost everywhere.
The latter now extends smoothly across it, so the two functions
define the same distribution there and $\mu$ has no atom at the
puncture.  Thus every alleged point of $Z$ is removable and has
zero atomic mass.  The preceding exponential-integrability and
Morrey estimates apply across it as well, and the original sequence
converges on all compact subsets of $\Delta$.  Since this applies to
every sequence, the family of centers is relatively compact.

If complete images could escape every compact, choose an escaping
interior point on each disc and precompose with a disc automorphism
sending zero to that point.  Both the boundary compact and the area
bound are unchanged, contradicting the center conclusion.  This
proves the first assertion.

For the cotangent assertion take global holomorphic fields
$X_1,\ldots,X_m$ spanning $T^{1,0}M$, as supplied by
Proposition~\ref{prop:initial-data}.  The base images are compact by
the first part.  The functions $\xi(X_i)$ along each disc are
bounded by their boundary values.  Over the resulting compact base
the evaluation map $\xi\mapsto(\xi(X_i))_i$ has a uniform positive
least singular value.  Its bounded values therefore bound the
cotangent vectors and put the images in a compact subset of $Y$.
Local holomorphic normal-family compactness proves the stated
subsequence conclusion.
\end{proof}

\subsection{Boundary factorization and holomorphic frames}

We now construct frames normalized on the boundary of each disc.
Their symplectic coefficients and their values at a fixed compact set
of centers must remain controlled when the disc area increases.
Fix the injective immersion
\begin{equation}\label{disc:immersion-y}
 G:Y\longrightarrow\C^q,\qquad
 G(x,\xi)=\big(F(x),\xi(X_1(x)),\ldots,\xi(X_m(x))\big).
\end{equation}
It is injective because its base entries separate base points and
the fiber evaluations separate covectors.  It is immersive because
$dF$ detects base tangent vectors and the evaluations detect vertical
ones.  There is also a finite family of global holomorphic fields
spanning $TY$: take the cotangent lifts
\[
 \widetilde X_i=\sum_aX_i^a\partial_{x_a}
       -\sum_{a,b}\xi_b(\partial_{x_a}X_i^b)\partial_{\xi_a}
\]
and the vertical translations by $dF_j$.
The former span after projection to $TM$, and the latter span the
vertical tangent.  Smooth coefficients on any fixed compact subset
of $Y$ can be represented by smooth functions of $G$ there, by local
embedded coordinate neighborhoods and a finite partition of unity.

For a vector or matrix function $a$ on the circle, write
\begin{equation}\label{disc:half-sobolev}
 D(a)^2=\sum_{k\in\mathbb Z}|k|\,\|a_k\|_F^2
 =\frac1{(2\pi)^2}\int_0^{2\pi}\!\int_0^{2\pi}
 \frac{\|a(e^{i(t+\theta)})-a(e^{it})\|_F^2}
 {|e^{i\theta}-1|^2}\,dt\,d\theta,
\end{equation}
where $a_k$ are Fourier coefficients and $\|\cdot\|_F$ is the
Frobenius norm.  For holomorphic $a$, this is
$\pi^{-1}\int_\Delta\|a'(z)\|_F^2\,dA$.

\begin{lemma}[Disc frames and Gauss area]\label{lem:disc-frames}
Let $f$ be a holomorphic disc in $Y$ extending past the unit circle.
There is a holomorphic frame $E_f$ of $f^*TY$, smooth on the closed
disc, for which $U=dG\,E_f$ satisfies
\begin{equation}\label{disc:frame-unitary}
 U^*U=I\quad\hbox{on }\partial\Delta,\qquad
 \sup_\Delta\|U\|_{\mathrm{op}}\le1.
\end{equation}
Let $\mathfrak g:Y\to\Gr(2n,q)$ send $y$ to $dG(T_yY)$, and
normalize the Pl\"ucker form by
$\omega_{\Gr}=i\partial\bar\partial\log\det(V^*V)$ for a local
holomorphic frame $V$ of its tautological plane.  Then
\begin{equation}\label{disc:gauss-area}
 D(U)^2=\frac1{2\pi}\int_\Delta
                          (\mathfrak g\circ f)^*\omega_{\Gr}.
\end{equation}
If $f(\partial\Delta)\subset K$ for a fixed compact $K\subset Y$,
the holomorphic component matrix $J_{E_f}$ of $\mathcal J$ in this
frame is bounded throughout $\Delta$ by a constant depending only
on $K$ and $G$.  If also $f(0)$ ranges in a compact subset $K_0$
of a fixed coordinate chart, the frame matrix at zero and its
inverse are bounded by constants depending only on $K$, $K_0$, $G$
and that chart.  These constants do not depend on the disc area or
on an interior compact containing its image.
\end{lemma}

We will first trivialize $f^*TY$ and then normalize the boundary Gram
matrix of that frame.  For the normalization we use the smooth
Wiener--Masani factorization \cite{WienerMasani1957} in the Hardy-space
form described by Berndtsson--Rosay
\cite[Theorem~2.1 and Section~6]{BerndtssonRosay}.
We include its construction and the boundary half-Sobolev estimates
that will control the later deformation error.

\begin{lemma}[Boundary factorization]\label{lem:boundary-factorization}
Let $L$ be an $r\times r$ smooth Hermitian matrix on the circle with
$mI\le L\le MI$, where $0<m\le M<\infty$.  There is a holomorphic
invertible matrix $H$ on $\Delta$, smooth with smooth inverse on
$\overline\Delta$, such that $L=H^*H$ on the circle.  It satisfies
\begin{align}
 \|H\|_{\infty,\mathrm{op}}&\le\sqrt M,&
 \|H^{-1}\|_{\infty,\mathrm{op}}&\le m^{-1/2},\label{disc:factor-sup}\\
 D(H)^2&\le\frac{M}{2m^2}D(L)^2,&
 D(H^{-1})^2&\le\frac1{2m^3}D(L)^2.
 \label{disc:factor-seminorm}
\end{align}
In particular the seminorm constants depend only on the boundary
spectral bounds, independently of higher derivatives of $L$.
\end{lemma}

\begin{proof}
Equip the vector Hardy space with
$\|v\|_L^2=\langle v^*Lv\rangle$, where brackets denote normalized
circle average.  This norm is equivalent to the ordinary Hardy norm.
The shift $v\mapsto zv$ is an isometry, its range has codimension
$r$, and $\bigcap_{k\ge0}z^kH^2=\{0\}$.  An orthonormal basis
$p_1,\ldots,p_r$ of the orthogonal complement of the range, followed
by all its successive shifts, is therefore an orthonormal basis of
the weighted space.  Indeed the orthogonal remainder after the first
$k$ shift levels is $z^kH^2$, whose intersection is zero.

Let $P$ be the holomorphic matrix with columns $p_i$.  Orthogonality
of all their shifts says that every Fourier coefficient of
$P^*LP-I$ vanishes, so $P^*LP=I$ almost everywhere on the circle.
Thus $P$ is bounded.  Multiplication by $P$ maps the ordinary Hardy
space isometrically onto the weighted space by the basis property.
Apply this to constant target columns to obtain a holomorphic matrix
$Q$ of Hardy functions with $PQ=I$.  Consequently $P$ is invertible
everywhere, $Q=P^{-1}$, and the boundary identity gives $L=Q^*Q$.
It also bounds both $P$ and $Q$ in $H^\infty$ by their boundary
values.  Taking $H=Q$ proves \eqref{disc:factor-sup}.

For the seminorm estimate put
$L_\theta(z)=L(e^{i\theta}z)$ and
$A_\theta(z)=H(e^{i\theta}z)H(z)^{-1}$.  This holomorphic matrix
has $A_\theta(0)=I$.  Its mean is therefore $I$, and
\[
 E_\theta:=\langle\|A_\theta-I\|_F^2\rangle
       =\langle\tr(L_\theta L^{-1}-I)\rangle.
\]
Changing variables in the expression for $E_{-\theta}$ gives
\begin{align*}
 E_\theta+E_{-\theta}
 &=\left\langle\tr\big((L_\theta-L)L^{-1}
                      (L_\theta-L)L_\theta^{-1}\big)\right\rangle\\
 &\le m^{-2}\langle\|L_\theta-L\|_F^2\rangle.
\end{align*}
The identities
\[
 H_\theta-H=(A_\theta-I)H,\qquad
 H_\theta^{-1}-H^{-1}=H_\theta^{-1}(I-A_\theta)
\]
bound their squared differences by $ME_\theta$ and
$m^{-1}E_\theta$, respectively.  Integrating against the symmetric
kernel in \eqref{disc:half-sobolev} proves
\eqref{disc:factor-seminorm}.

It remains to justify smooth boundary regularity.  The same difference
inequality, divided by $\theta^2$, first gives one full $L^2$
tangential derivative of $H$ and $H^{-1}$.  With prime denoting that
derivative, set $B=H'H^{-1}$; this is the boundary value of
$izH_zH^{-1}$ and has zero constant Fourier coefficient.  Differentiating
$L=H^*H$ gives
\[
 H^{-*}L'H^{-1}=B+B^*.
\]
The strictly positive Fourier projection recovers $B$.  The
one-dimensional Sobolev algebra property now bootstraps: if
$H,H^{-1}\in W^{k,2}$, the displayed expression is in $W^{k,2}$,
hence so is $B$, and $H'=BH$ gives $H\in W^{k+1,2}$.
The inverse derivative formula gives the same for $H^{-1}$.
For this one-dimensional algebra assertion, Cauchy--Schwarz on
Fourier coefficients and $\sum_k(1+k^2)^{-1}<\infty$ give
$W^{1,2}(S^1)\subset L^\infty(S^1)$.  In each Leibniz term of total
order at most $k\ge1$, place a derivative of order at most $k-1$
in $L^\infty$ and the other factor in $L^2$; this gives the
$W^{k,2}$ product estimate.  The positive-frequency projection has
norm at most one in every $W^{k,2}$.  Iterating the preceding
identities therefore gives all Sobolev orders, and the same Fourier
estimate for derivatives proves smoothness on the closure.
\end{proof}

\begin{proof}[Proof of Lemma~\ref{lem:disc-frames}]
We first construct a preliminary frame on a neighborhood of the closed
disc.  The bundle $E=f^*TY$ is embedded by $dG$ in a trivial bundle
and is generated by the pullbacks of the finite global spanning
fields.  Start on a disc of radius greater than one and choose a
generator not identically zero.  On a slightly smaller disc, divide
it by the finitely many common zero factors of its ambient
components, using at each zero the minimum component vanishing
order.  The resulting section $e$ is holomorphic, remains in $E$ by
continuity, and is nowhere zero.

Write its ambient components as $e_1,\ldots,e_q$ and put
$\ell_j=\overline{e_j}/|e|^2$.  Thus $\sum_j e_j\ell_j=1$.
The smooth $(0,1)$-forms
\[
 b_{jk}=\ell_j\dbar\ell_k-\ell_k\dbar\ell_j
\]
are antisymmetric and satisfy
$\sum_j e_jb_{jk}=\dbar\ell_k$.
Solve $\dbar c_{jk}=b_{jk}$ with $c_{jk}=-c_{kj}$ on a further
slightly smaller disc.  This scalar step follows directly by
multiplying the coefficient of each $(0,1)$-form by one common smooth
compactly supported cutoff equal to one on that disc and convolving
on $\C$ with $1/(\pi z)$;
the distributional identity
$\dbar(1/(\pi z))=\delta_0$ gives a smooth solution for the smooth
source.  It does not use a frame of $E$.
The functions
\[
 a_k=\ell_k-\sum_j e_jc_{jk}
\]
are holomorphic and obey $\sum_k e_ka_k=1$.
The ambient row $a$ restricts to a holomorphic dual of $e$ on $E$,
so $E=\C e\oplus\ker a$.  If $s_i$ are the current generators,
then $s_i-e\,a(s_i)$ generate $\ker a$.
Thus the kernel is again an ambiently embedded, finitely generated
holomorphic bundle, now of rank one less.  Repeat the construction
on a finite nested sequence of discs whose radii remain greater
than one.  Induction gives a holomorphic frame of $f^*TY$ on a
neighborhood of $\overline\Delta$.

Choose this preliminary frame and apply
Lemma~\ref{lem:boundary-factorization} to its $dG$ Gram matrix on
the circle.  Multiplication of the frame by $H^{-1}$ yields $E_f$.
The first identity in \eqref{disc:frame-unitary} follows immediately.
For each constant column $v$, the maximum principle applied to the
holomorphic vector $Uv$ gives $|Uv|\le|v|$, proving the second.

Put $Q=U^*U$.  Since $Q=I$ on the circle,
\[
 \partial_r\tr Q=\tr(\partial_rQ)
                =\partial_r\log\det Q
 \quad\hbox{on }\partial\Delta.
\]
Green's formula and holomorphicity give
\[
 4\pi D(U)^2
 =\int_\Delta\Delta_{\mathrm{eucl}}\tr Q\,dA
 =\int_\Delta\Delta_{\mathrm{eucl}}\log\det Q\,dA
 =2\int_\Delta(\mathfrak g\circ f)^*\omega_{\Gr},
\]
which proves \eqref{disc:gauss-area}.

On the boundary, $E_f$ is orthonormal for the metric $G^*g_{\rm eucl}$.
Thus the components of $\mathcal J$ there are bounded by its norm
in that metric on $K$.  They are holomorphic functions in the
frame, so the maximum principle gives the same bound inside.
For the last assertion, write $E_f=\partial S$ in the fixed chart
at the center, and let $J$ be the invertible coordinate matrix of
$\mathcal J$.  Then
\[
 J_{E_f}=S^{-1}JS^{-\mathsf T},\qquad
 S^{-1}=J_{E_f}S^{\mathsf T}J^{-1}.
\]
Let $\gamma>0$ be the infimum on $K_0$ of the least singular value
of $dG$ in the coordinate Euclidean metric.  By
\eqref{disc:frame-unitary}, $\|S(0)\|_{\mathrm{op}}\le\gamma^{-1}$.
The preceding identity gives
\[
 \|S(0)^{-1}\|_{\mathrm{op}}
 \le \sup_\Delta\|J_{E_f}\|_{\mathrm{op}}\,
                \gamma^{-1}\sup_{K_0}\|J^{-1}\|_{\mathrm{op}}.
\]
Every quantity on the right has the stated dependence.
\end{proof}
